\documentclass[11pt]{article}
\usepackage[a4paper,margin=1in]{geometry}
\usepackage[T1]{fontenc}
\usepackage[utf8]{inputenc}
\usepackage{lmodern}
\usepackage{booktabs}
\usepackage{array}
\usepackage{tabularx}
\usepackage{longtable}

\setlength{\parindent}{0pt}
\setlength{\parskip}{0.7em}
\renewcommand{\arraystretch}{1.15}

\begin{document}

\begin{center}
{\LARGE \textbf{TransFit Model and Parameter Reference}}
\end{center}

\section*{Recommended Public Model Keys}

\begin{longtable}{>{\ttfamily}p{0.18\textwidth} p{0.22\textwidth} p{0.52\textwidth}}
\toprule
\textnormal{\textbf{Model key}} & \textbf{Physical meaning} & \textbf{Parameter order} \\
\midrule
\endhead
\texttt{nickel} &
Ni-powered model &
\texttt{(M\_ej, v\_ej, M\_Ni, x\_Ni, kappa, kappa\_gamma, T\_floor)} \\

\texttt{sc\_ni} &
Shock-cooling + Ni &
\texttt{(M\_ej, v\_ej, E\_Th\_in, M\_Ni, R\_0, x\_Ni, kappa, kappa\_gamma, T\_floor)} \\

\texttt{magnetar} &
Pure magnetar &
\texttt{(M\_ej, v\_ej, P\_ms, B14, kappa, kappa\_gamma, T\_floor)} \\

\texttt{magnetar\_ni} &
Magnetar + Ni &
\texttt{(M\_ej, v\_ej, P\_ms, B14, M\_Ni, kappa, kappa\_gamma, T\_floor)} \\

\texttt{sc\_magnetar} &
Shock-cooling + magnetar &
\texttt{(M\_ej, v\_ej, E\_Th\_in, P\_ms, B14, R\_0, kappa, kappa\_gamma, T\_floor)} \\
\bottomrule
\end{longtable}

Compatibility aliases are still accepted internally, but the names above are the recommended public API.

\section*{Parameter Glossary}

\begin{tabularx}{\textwidth}{>{\ttfamily}p{0.18\textwidth} X >{\ttfamily\raggedright\arraybackslash}p{0.20\textwidth}}
\toprule
\textnormal{\textbf{Parameter}} & \textbf{Meaning} & \textbf{Unit} \\
\midrule
\texttt{M\_ej} & Ejecta mass & \texttt{M\_sun} \\
\texttt{v\_ej} & Characteristic ejecta velocity & \texttt{1e9 cm s\^{}-1} \\
\texttt{M\_Ni} & Nickel-56 mass & \texttt{M\_sun} \\
\texttt{E\_Th\_in} & Initial thermal energy scale & \texttt{1e49 erg} \\
\texttt{R\_0} & Initial outer-radius scale & \texttt{R\_sun} \\
\texttt{x\_Ni} & Heating radius fraction & dimensionless \texttt{[0,1]} \\
\texttt{kappa} & Optical opacity & \texttt{cm\^{}2 g\^{}-1} \\
\texttt{kappa\_gamma} & Gamma-ray opacity & \texttt{cm\^{}2 g\^{}-1} \\
\texttt{P\_ms} & Magnetar initial spin period & \texttt{ms} \\
\texttt{B14} & Magnetar magnetic field scale & \texttt{1e14 G} \\
\texttt{T\_floor} & Temperature floor for photosphere & \texttt{K} \\
\texttt{t\_shift} & Time offset between model and observed timeline & \texttt{day} \\
\bottomrule
\end{tabularx}

\section*{\texttt{t\_shift} Convention}

In fitting mode, the likelihood is evaluated as
\[
\mathrm{model}(t_{\mathrm{obs}} + t_{\mathrm{shift}}).
\]

Therefore:
\begin{itemize}
  \item positive \texttt{t\_shift} shifts the model curve to earlier observed times;
  \item if you do not want to fit \texttt{t\_shift}, fix it with \texttt{fixed = \{"t\_shift": 0.0\}}.
\end{itemize}

\end{document}
